\section{Scope and open problems}\label{sec:scope}

\subsection{Pointwise T2}\label{sec:t2open}

F3's frozen scope sentence is the boundary and I repeat it verbatim: F3
measures the viscous enstrophy flux of the $\omega_1$ channel as an
integrated functional over $\nu \in [10^{-8}, 10^{-5}]$ within the
stored collapse window; it does not resolve the pointwise T2 question,
which remains open
(\texttt{ext\_inertia\_t2/out/flux3/F3\_RESULT.md}, section 6). Nothing
in this paper constrains
$\omega_{1,\nu} - \omega_{1,\mathrm{inv}} = \nu\,\delta_1 + o(\nu)$
pointwise, least of all at the corner ring; nothing is in sup-norm;
rearrangements preserving the gradient integral are invisible to $Z$;
window-local violations can cancel inside a window, mitigated but not
removed by the nine-window report. Section~4 explains why the pointwise
question is hard here rather than merely unfinished: the sealed record
has viscosity acting through the collapse frame itself, and a
frame-quotiented statistic deleted 71.8 to 99\% of its own prediction
(the root-cause table, Section~4). A pointwise instrument for this
collapse therefore must not quotient the frame, and no such instrument
with a theorem-backed null is in hand. The one concrete function-space
lead comes from the corner treatment of Chen and Hou, who carry the
corner with singular weights and explicit corrections rather than
excision \cite{ChenHou}: a pointwise corner response coefficient would
need that function-space change, not grid refinement.

\subsection{The sub-leading exponent}\label{sec:subleading}

The joint license's floor, corrections vanishing no slower than
$\nu^{0.4}$, is a band edge, not a measurement. The F3 bulk correction
fit returns $q = 0.438$ (Simpson, window $[200, 400)$) and $q = 0.425$
(trapezoid, window $[1200, \Omega_*)$, jackknife spread $0.088$)
against the frozen band $[0.4, 1.0]$: both sit at the edge, so the
exponent is unresolved
(\texttt{ext\_inertia\_t2/out/flux3/f3\_ladder.json},
\texttt{family\_A}; \texttt{F3\_RESULT.md}, section 9). The F2 swirl
correction exponent at fixed time is $0.7130$, drop-one jackknife
spread $0.029$, profile $[0.710, 0.716]$, and it lives at
$\nu \in [10^{-5}, 10^{-3}]$
(\texttt{ext\_inertia\_t2/out/flux/FLUX\_RESULT.md}, section 4.2).
Different channel, different pairing, different $\nu$ range: I do not
reconcile the two numbers and this archive cannot. Resolving $q$ is a
floor problem; see the archive limits below.

\subsection{R2: azimuthal stability, parked, with the certified route}\label{sec:r2parked}

Whether the corner ring survives azimuthal symmetry breaking is not
addressed by any measurement in this paper. The question is parked with
a certified route on disk. The Tier-0 pre-registration
(\texttt{ext\_inertia\_t2/R2\_PLAN/TIER0\_SPEC.md}) froze the route: an
m-mode operator about the certified profile, azimuthal coupling in a
polynomial normal form
$L_m = L_0 + \tilde{k} B_1 + \tilde{k}^2 B_2$ at
$\tilde{k} = m\,\ell(t)$, with the certified margin
$m_{\mathrm{cert}} = 1/(2\lambda_{\max}(P))$ from the Lyapunov
certificate. The execution
(\texttt{ext\_inertia\_t2/R2\_PLAN/TIER0\_RESULT.md}) returned T0c,
VOID by the spec's own stop rule, and the void has content. Certified
and re-derived from $P$:
$m_{\mathrm{cert}} = 8.8232737664\times10^{-6}$, and
$\mathrm{Re}\,\lambda \le -8.8233\times10^{-6}$ for every eigenvalue of
the reduced generator, no eigensolve used. Measured obstructions: the
2D certificate cannot host the azimuthal couplings, since $82.71\%$ of
its nodes sit at cylindrical radius $R \le 0$ and those rows carry
$91.27\%$ of $\mathrm{tr}(P)$, with the $1/R$ pole interior to the
discretization; and my polynomial ansatz was wrong twice over, because
the dominant azimuthal channel is the swirl advection entering at
$\kappa = m\,(t_s - t)$, which exceeds $\tilde{k}$ by
$\ell^{-0.658}$ at $c_l = 2.9206$, and because pressure slaving makes
the $m$-dependence operator-rational rather than polynomial. The honest
normal form is two-parameter, $(\kappa, \tilde{k})$, and the certified
band depth for a fixed $m$ scales as $m^{-1}$ (swirl channel), not
$m^{-0.342}$. The corner circle, the one locus where the azimuthal
coefficients are unambiguous ($R = 1$ exactly at $\rho = 0$), is
exactly the ring the certificate deletes. The closing route is
itemized in the result file, first item a 3D corner similarity system;
the fixed-$\tilde{k}$ ladder remains the right experiment. Two
disciplines bind any future R2 run. First, perturbation growth is never
a falsifier: blowup solutions of this type are hydrodynamically
unstable by theorem \cite{VasseurVishik2020}, so discriminators must be
structural. Second, no unconditional linear stability of the profile is
claimed here or anywhere in this paper: the certificate is
corner-clamped, in the P-norm, with measured transient growth
$89.7 \le \sup_t \|e^{tL}\| \le 5807$ (Section~2), and that caveat
travels with every mention of the certificate.

\subsection{Archive limits, and what a purpose-built archive would buy}

The deep end of this archive is a floor property, not an estimator
property. In the swirl channel the fixed-time deviation falls to the
inviscid instrument floor ($1.80\times10^{-10}$ on the D and T grids)
below $\nu \approx 3\times10^{-6}$; the deepest rung that ever clears
the deviation gate is $3\times10^{-6}$, on D2048 at
$t \ge 2.85\times10^{-3}$, so the F2 exponent lives at
$\nu \in [10^{-5}, 10^{-3}]$ and nothing below $10^{-6}$ is decidable
there (\texttt{ext\_inertia\_t2/out/flux/FLUX\_RESULT.md}, sections 6.2
and R1). In the $\omega_1$ channel the $10^{-8}$ rung's deepest window
is floor-unrealized under the primary scheme (cross-realization ratio
$12.7$ on roundoff-scale floors) and is decided by the trapezoid
control in class agreement; the grid sigma term, $0.068$, forced by the
absent \texttt{snap\_D2048\_1e-6} cell, is the largest single error
term on every T rung; the floors are single-realization on two of the
three grids (in the swirl lane the T and H floors inherit the D-grid
replicate fraction, $1.20\%$,
\texttt{ext\_inertia\_t2/out/flux/FLUX\_RESULT.md}, section 2.5; in the
$\omega_1$ lane no inviscid field replicate exists at all and the
floor-realization spread across independent executions is unmeasured,
substituted by the 5 to 7\% cross-grid floor transfer,
\texttt{F3\_RESULT.md}, exposure E-1); and the thin deep bin
$[\Omega_*, 1900)$ holds five intervals at $10^{-8}$
(\texttt{ext\_inertia\_t2/out/flux3/F3\_RESULT.md}, sections 5 and 9,
exposures E-1..E-3). A purpose-built archive would need: an inviscid
replicate pair on the T grid; snapshot storage, not streams, for every
rung entering a fit; denser write cadence in the corner window
$\Omega \in [1200, 1900)$; and the $\nu = 10^{-6}$ same-depth D/T cell.
That list buys floors and cross-checks. It creates no new physics
claims, and I make none from it.

\subsection{What is not claimed}

Stated explicitly, each with the reach that bounds it.

\begin{enumerate}
\item Nothing about generic Navier-Stokes blowup. Every result in this
paper is about one certified profile
($\alpha = -0.34240 \pm 3\times10^{-5}$, Section~2), one geometry, one
boundary family.
\item Nothing about no-slip walls beyond the wall-swap probe of
Section~\ref{sec:wall}. The no-slip ladder was a boundary-condition
experiment on the clamp hypothesis, thirteen completed short runs,
formally W3;
the no-slip half-space regularity theory is prior art
\cite{GigaMiura2011, BarkerPrange2020} and this campaign adds nothing
to it. A direct no-slip comparison of the direction-regularity exponent
was attempted on the stored no-slip family (18 August 2026, pre-registered,
census-gated) and refused itself: the impulsive-start Rayleigh sheet attains
the vorticity supremum on every gated snapshot of all fourteen runs, at the
quarter-period rings $z = L/4,\,3L/4$ rather than at the corner, so the
certified observable reads a different object under no-slip walls and no
cross-boundary-condition comparison is licensed from that data.
\item Nothing pointwise, at any $\nu$. T2 is open
(Section~\ref{sec:t2open}).
\item No single scalar sub-leading exponent. The matching-convention
fork fires in the swirl channel (VOID-5 / V-4, Section~5) and the
bulk-channel $q$ sits at its band edge
(Section~\ref{sec:subleading}).
\item Nothing below $\nu = 10^{-8}$, and nothing deeper than the stored
window ($t \le 2.985\times10^{-3}$, $\Omega \lesssim 1780$). No number
in this paper is an extrapolation beyond its measured range.
\item No corner-mechanism attribution. The corner share never reaches
its gate (maximum $0.412$ against $0.5$, Section~\ref{sec:wall}), and
corner language is licensed in neither flux lane.
\item No unconditional linear stability of the profile, and no claim
about azimuthal stability: R2 is parked
(Section~\ref{sec:r2parked}), and the certificate's transient-growth
caveat stands.
\item No anomalous-dissipation statement outside this window and
geometry. The zero-anomaly verdict of Section~3 concerns this collapse
through $\Omega_*$; it says nothing about Onsager-class dissipation in
turbulence \cite{Onsager1949, DuchonRobert2000}.
\end{enumerate}
